% !TeX program = xelatex
\documentclass{jos}
\makeatletter
\renewcommand{\jos@topmatter}{\noindent 天津大学 NLP 平时作业\hfill 第8队个人阅读模块\par}
\renewcommand{\jos@citationbox}[2]{}
\makeatother
\fancyhead[RO]{\joskaiti Semantic XPath 阅读报告\quad 第8队}
\setlength{\headheight}{16pt}
\begin{document}
\Year{2026}\BeginPage{1}\CLC{TP18}
\title{Semantic XPath结构化记忆访问方法阅读报告}{Semantic XPath阅读报告}
\entitle{A Reading Report on Semantic XPath for Structured Conversational Memory Access}{Semantic XPath Reading Report}
\author{oioioi-521}{oioioi-521}{}
\AuthorMark{oioioi-521}\address{天津大学 NLP 第8队}{NLP Team 8, Tianjin University}
\abstract{Semantic XPath将对话任务记忆组织为树，结合路径、位置选择与连续语义评分访问相关节点和子树。本报告梳理数据模型、查询执行与证据聚合，核对作者单轮、多轮与LoCoMo结果，区分任务通过率、回答得分和token口径。该方法改善相对扁平检索的效果及相对全量上下文的输入规模，但准确性取决于模型和任务。本报告进一步讨论层级依赖、证据完整性、元数据过滤基线及与本组课设的联系。}
\keywords{结构化记忆；对话智能体；Semantic XPath；语义检索；论文阅读}
\maketitle
\section{Semantic XPath}
Semantic XPath\upcite{liu-etal-2026-semantic}研究结构化对话记忆的访问。它结合路径、位置条件和连续语义评分，从已有树形记忆中选择相关节点及子树，支持下游回答和状态更新。

\subsection{研究问题与方法定位}

长期任务导向对话需要反复查看、修改和回溯同一个任务对象。旅行计划中的一项活动不仅包含文本，还属于特定日期、行程与历史版本。全量上下文保留信息较多，但输入量随历史增长；扁平检索降低输入量，却可能取回语义相似而所属范围错误的条目。比如“在会议安排最密集的一天增加咖啡休息”，需要比较各天的活动，而不仅是找一条包含会议内容的文本。

查询负责定位访问范围，插入、删除及版本维护由下游步骤完成。

\subsection{树形数据模型与查询语法}

记忆表示为有根树 $T=(V,E,r)$。非根节点只有一个父节点，节点具有类型和文本属性；模式 $\Sigma$ 由记忆数据已有的节点类型、属性和父子关系决定，而非要求人工预先设计一个通用模式。旅行示例采用
\[
\text{Itinerary}\rightarrow\text{Version}\rightarrow\text{Day}\rightarrow\text{POI},
\]
分别表示计划、版本、日期和活动；层级主要是组成关系，不只是对下层文本做有损摘要。

查询是步骤序列，每一步写为 $S=A\,N\,[R]\,[P]$：路径与类型必需，位置选择器和语义条件可选。
\begin{table}[H]
\centering
\caption{Semantic XPath 查询步骤的组成}
\label{tab:sxpath-grammar}
\begin{tabularx}{\linewidth}{@{}p{20mm}p{43mm}X@{}}
\toprule
组成 & 示例写法 & 作用 \\
\midrule
路径 $A$ & \texttt{/}、\texttt{//} & 访问直接子节点或全部后代 \\
类型 $N$ & \texttt{Day}、\texttt{POI}、\texttt{*} & 按节点类型筛选或匹配任意类型 \\
位置 $R$ & \texttt{[i]}、\texttt{[-i]}、\texttt{[i:j]} & 按顺序、倒序或区间选取候选 \\
语义 $P$ & \texttt{[node\textasciitilde="workshop"]} & 计算候选与语义条件的相关程度 \\
\bottomrule
\end{tabularx}
\end{table}

例如，\texttt{//Day[3]/POI[node\textasciitilde="workshop"]}先选第三个匹配日期，再对该日期下的活动评分。存在多个行程或版本时，查询需先限定目标范围。语义条件$1-P$对原条件评分取反，提高低分候选的优先级。

\subsection{语义评分与带权执行}

语义相关度定义为 $\operatorname{Rel}(u,P)\in[0,1]$。局部评分依据一个属性或完整节点表示；聚合评分先在当前节点下执行内部结构查询，再用平均、最小、最大或几何平均聚合证据；多个条件还可用乘积、平均及取反组合。

\begin{quote}\small\ttfamily
//Day[avg(/POI[node\textasciitilde="conference"])]
\end{quote}
这个查询根据活动证据评价日期。示例中第二天的三个活动分数为$0.603$、$0.482$、$0.608$，平均约$0.565$，在候选日期中得分最高。该值衡量活动与会议主题的平均相关度。

执行器从根节点及权重1出发，维护带权集合 $W\subseteq V\times[0,1]$，按步骤进行路径扩展、类型筛选、位置选择与语义评价。语义步骤更新为
\begin{equation}
w_{\mathrm{new}}=w_{\mathrm{old}}\operatorname{Rel}(u,P).
\end{equation}
聚合条件递归访问下层证据。阈值、返回数量和空集处理由具体任务设置，查询质量则可从目标范围、节点选择和证据完整性三个方面检查。

\subsection{范围约束示例}

以行程调整为例，假设当前版本的第2天安排海滨公园，第5天安排山地徒步，第7天安排植物园和海港游船。用户要求找出“第7天下雨时需要调整的户外活动”。纯语义排序可能混入其他日期的活动；结构化访问先限定第7天，再评价活动语义，目标为植物园与游船。

简单日期约束也可由带元数据过滤的扁平RAG完成。比较查询机制时，可将其作为基线，进一步考察嵌套范围、下层证据聚合等结构访问需求。

\subsection{实验设置与指标口径}

作者评估旅行行程、待办清单和家庭配餐三个自建场景，每场景20个单轮对话、5个多轮对话；另在 LoCoMo 中使用原生的 Memory--Conversation--Session--Turn 层级。回答模型为 GPT-5 mini 与 Gemini 3 Flash。相似度评分器（Sem）和 Flat RAG 使用 Qwen3-Embedding-8B；蕴含评分器（Entail）使用 \texttt{facebook/bart-large-mnli}。

自建场景使用任务通过率，要求检索到正确结构化数据且输出满足任务约束；LoCoMo使用模型裁判评定的回答得分。LoCoMo的平均输入tokens包含XPath查询生成输入，自建场景则报告平均token用量。成本比较还需计入嵌入、蕴含评分和查询执行，并在同一模型内对齐分词与统计口径。

\begin{table}[H]
\centering
\caption{三个自建场景的平均单轮结果}
\label{tab:sxpath-single}
\begin{tabular}{@{}llrr@{}}
\toprule
回答模型 & 方法 & 通过率（\%） & 平均tokens \\
\midrule
GPT-5 mini & In-context & 70.0 & 26\,216 \\
 & Flat RAG & 30.0 & 3\,250 \\
 & XPath Sem & 66.3 & 6\,166 \\
 & XPath Entail & 83.0 & 5\,691 \\
\midrule
Gemini 3 Flash & In-context & 79.0 & 25\,387 \\
 & Flat RAG & 37.7 & 3\,361 \\
 & XPath Sem & 64.3 & 6\,043 \\
 & XPath Entail & 70.7 & 5\,830 \\
\bottomrule
\end{tabular}
\end{table}

单轮GPT-5 mini配置中，Entail相对Flat RAG提高53.0个百分点，相对增幅约176.7\%；token用量约为全量上下文的21.7\%。Gemini配置中，全量上下文的79.0\%高于Entail的70.7\%，评分机制的收益随模型变化。

\begin{table}[H]
\centering
\caption{三个自建场景的平均多轮结果}
\label{tab:sxpath-multi}
\begin{tabular}{@{}llrr@{}}
\toprule
回答模型 & 方法 & 通过率（\%） & 平均tokens \\
\midrule
GPT-5 mini & In-context & 80.0 & 69\,097 \\
 & Flat RAG & 26.7 & 3\,327 \\
 & XPath Sem & 70.0 & 6\,294 \\
 & XPath Entail & 83.3 & 5\,850 \\
\midrule
Gemini 3 Flash & In-context & 56.7 & 76\,603 \\
 & Flat RAG & 26.7 & 3\,240 \\
 & XPath Sem & 76.7 & 6\,092 \\
 & XPath Entail & 80.0 & 5\,993 \\
\bottomrule
\end{tabular}
\end{table}

多轮Entail在两种模型下取得最高平均通过率，token用量约为全量上下文的8.5\%与7.8\%。Sem与Entail的差异说明评分器选择会影响任务效果。单轮与多轮实验使用不同请求，分别反映两类任务设置。

\begin{table}[H]
\centering
\caption{LoCoMo结果（XPath输入含查询生成）}
\label{tab:sxpath-locomo}
\begin{tabular}{@{}llrr@{}}
\toprule
回答模型 & 方法 & 回答得分 & 平均输入tokens \\
\midrule
GPT-5 mini & In-context & 74.55 & 24\,857 \\
 & Flat RAG & 35.71 & 1\,341 \\
 & Semantic XPath & 65.19 & 3\,032 \\
\midrule
Gemini 3 Flash & In-context & 87.01 & 27\,499 \\
 & Flat RAG & 41.49 & 1\,996 \\
 & Semantic XPath & 73.25 & 3\,870 \\
\bottomrule
\end{tabular}
\end{table}

LoCoMo上XPath的总体得分处于两种基线之间，输入分别约为全量上下文的12.2\%与14.1\%。这说明已有会话层级可用于访问，但准确性仍有损失。例如 GPT-5 mini 多跳题得分为43.26，低于全量上下文58.87，证据整合仍有差距。作者的失败分析将错误分为查询生成、检索和回答生成，认为检索是主要剩余来源，包括无命中、证据不全及选错节点或发言。

\subsection{方法局限与课设启发}

结构范围与连续语义评分支持不同粒度的访问，路径与评分轨迹便于解释结果。方法依赖可靠层级，跨树关系、模糊版本和复杂多跳仍有困难。现有评测规模较小，重复实验、置信区间及本地评分开销也有待补充。

StructMem侧重写入时保存事件关联，Semantic XPath侧重读取时限定范围。课程设计可明确日期与版本约束，保留节点及来源标识，分别检查范围、证据完整性和回答错误，并以元数据过滤基线评估层级访问的贡献。

% 原文定位：PDF第1--2页，图1--2；PDF第3页，\S3.1；PDF第3--4页，\S3.2；PDF第10页；PDF第4页，\S3.3、图2；PDF第5页，\S4.1；PDF第6--7页，\S4.3、图8。

\bibliographystyle{jos}\bibliography{references}
\end{document}
